% !TEX program = xelatex
\documentclass[11pt]{article}
\usepackage[a4paper,margin=24mm]{geometry}
\usepackage{amsmath,amssymb,amsthm,mathtools}
\usepackage{xcolor}
\usepackage[colorlinks=true,linkcolor=blue!50!black,urlcolor=blue!50!black,citecolor=blue!50!black]{hyperref}
\newcommand{\C}{\mathbb C}
\newcommand{\Z}{\mathbb Z}
\newcommand{\g}{\widehat{\mathfrak{sl}}_2}
\newcommand{\gt}{\widetilde{\mathfrak{sl}}_2}
\newcommand{\ad}{\operatorname{ad}}
\newcommand{\Span}{\operatorname{span}_{\C}}
\newcommand{\I}{\mathcal I}
\newcommand{\D}{\mathcal D}
\newcommand{\wt}{\operatorname{wt}}
\newtheorem{theorem}{Theorem}[section]
\newtheorem{lemma}[theorem]{Lemma}
\newtheorem{corollary}[theorem]{Corollary}
\theoremstyle{definition}
\newtheorem{remark}[theorem]{Remark}
\setlength{\emergencystretch}{2em}
\title{Boundary Localization and Adjoining\\the Degree Derivation}
\author{Supplementary research note}
\date{September 8, 2026}
\begin{document}
\maketitle

\section{Definitions, hypotheses, and scope}

All vector spaces and tensor products are over \(\C\). Write
\[
 \g=(\mathfrak{sl}_2\otimes\C[t,t^{-1}])\oplus\C K,
 \qquad U=U(\g)
\]
for the affine algebra without the degree derivation. For
\(x\in\{e,h,f\}\) and \(r\in\Z\), set \(x_r=x\otimes t^r\).
Our bracket conventions are
\begin{equation}\label{eq:brackets}
\begin{aligned}
 [e_r,f_s]&=h_{r+s}+r\delta_{r+s,0}K,
 &[h_r,h_s]&=2r\delta_{r+s,0}K,\\
 [h_r,e_s]&=2e_{r+s},
 &[h_r,f_s]&=-2f_{r+s},
\end{aligned}
\end{equation}
where \(\delta\) is the Kronecker delta, \(K\) is central, and
\([e_r,e_s]=[f_r,f_s]=0\). Adjoining the degree derivation gives
\[
 \gt=\g\rtimes\C d,\qquad [d,x_r]=rx_r,\qquad[d,K]=0,
 \qquad\widetilde U=U(\gt).
\]

Fix \(a,b\in\Z\) with \(a+b\ge0\), and put \(L=a+b+1\ge1\). Define
\[
 S_{a,b}=\Span\{K,h_i\ (i\ge0),e_i\ (i>a),f_j\ (j>b)\}.
\]
The bracket relations show that this is a Lie subalgebra and that
\[
 [S_{a,b},S_{a,b}]
 =\Span\{e_i\ (i>a),f_j\ (j>b),h_k\ (k\ge L+1)\}.
\]
Choose \(c,\lambda_0,\ldots,\lambda_L\in\C\), and set \(\lambda_i=0\)
for \(i>L\). These parameters define a character
\(\chi:S_{a,b}\to\C\) by
\[
 \chi(K)=c,\qquad\chi(h_i)=\lambda_i\quad(i\ge0),\qquad
 \chi(e_i)=\chi(f_j)=0,
\]
with the root indices restricted to those occurring in \(S_{a,b}\).
Set
\[
 M_{a,b}(\chi)=U\otimes_{U(S_{a,b})}\C_\chi,\qquad
 \C_\chi=\C1_\chi,\quad s1_\chi=\chi(s)1_\chi,\quad
 u=1\otimes1_\chi.
\]
In particular, \(h_0\) belongs to the inducing subalgebra. These conventions
agree with \cite[\S2.3]{FGXZ}. The main isomorphism requires
\(\lambda_L\ne0\); the central parameter \(c\) is arbitrary.

For \(F=f_n\), write
\[
 U_F=U[F^{-1}],\qquad
 \D_FM=U_F\otimes_U M,\qquad
 T_FM=\operatorname{coker}(M\longrightarrow\D_FM).
\]
The Ore localizations used here are justified in Section~4.
When \(F\) acts injectively, we identify \(M\) with its image in
\(\D_FM\) and write \(T_FM=\D_FM/M\). For \(\gt\)-modules, use
\(\widetilde U_F\) in the same definitions. An expression \(F^{-m}u+M\)
denotes a coset; it does \emph{not} assert that \(F^{-1}\) acts on the
quotient module. A module \(V\) is called \emph{smooth} if, for each
\(v\in V\), all sufficiently high modes \(e_r,h_r,f_r\) annihilate \(v\).

\paragraph{Two distinct questions.}
Replacing \(f_1\) by \(f_n\) may mean moving the inducing boundary along
with \(n\), or keeping the original module \(M_{-1,1}\) fixed. The first
interpretation admits a uniform isomorphism theorem; the second does not
always do so. Likewise, adjoining \(\C[d]\) requires a specified action
satisfying \([d,x_r]=rx_r\). An unshifted tensor-product action does not
in general satisfy this relation.

\section{Localization at an arbitrary boundary}

\begin{theorem}\label{thm:boundary}
Assume \(a+b\ge0\), \(L=a+b+1\), and \(\lambda_L\ne0\). Let
\(M=M_{a,b}(\chi)\) and \(F=f_b\). Define a character on \(S_{a+1,b-1}\) by
\[
 \chi^+(K)=c,\qquad
 \chi^+(h_i)=\lambda_i+2\delta_{i,0}\quad(i\ge0),
\]
with value zero on all root vectors in that subalgebra. Let \(v\) be the
canonical inducing vector of \(M_{a+1,b-1}(\chi^+)\). There is a unique
\(\g\)-module isomorphism
\begin{equation}\label{eq:boundary-iso}
 \Phi:M_{a+1,b-1}(\chi^+)\xrightarrow{\sim}T_{f_b}M_{a,b}(\chi),
 \qquad \Phi(v)=F^{-1}u+M.
\end{equation}
This holds for every \(c\in\C\), including the critical level \(c=-2\),
and also allows \(b\le0\).
\end{theorem}

\subsection{PBW bases and well-definedness}

Put \(Y=e_{a+1}\) and \(A=M_{a+1,b-1}(\chi^+)\). Let \(\mathcal X\)
consist of the finite ordered monomials
\[
 X=e(\alpha)h(\beta)f(\gamma),\qquad
 \alpha_i\le a,\quad\beta_i\le-1,\quad\gamma_i\le b-1.
\]
Here \(\alpha,\beta,\gamma\) are finite nondecreasing sequences,
\(e(\alpha)=e_{\alpha_1}\cdots e_{\alpha_{\ell(\alpha)}}\), and similarly
for \(h(\beta)\) and \(f(\gamma)\). Empty products are allowed, and
\(\ell\) denotes the number of factors. Define
\[
 \rho(X)=2\ell(\alpha)+\ell(\beta),\qquad
 w_m=F^{-m}u+M,\qquad B_{m,X}=F^{-m}Xu+M\quad(m\ge1).
\]
For \(M\), choose a PBW order with \(F\) first, followed by the
generators occurring in \(X\), and then a basis of \(S_{a,b}\). Thus
\(\{F^qXu:q\ge0,\ X\in\mathcal X\}\) is a basis of \(M\), so \(F\)
acts injectively. Clearing a common denominator and using this basis
shows that \(\{F^qXu:q\in\Z,\ X\in\mathcal X\}\) is a basis of
\(\D_FM\). Consequently,
\begin{equation}\label{eq:pbw-bases}
\begin{aligned}
 &\{B_{m,X}:m\ge1,\ X\in\mathcal X\}
     &&\text{is a basis of }T_FM,\\
 &\{XY^pv:p\ge0,\ X\in\mathcal X\}
     &&\text{is a basis of }A.
\end{aligned}
\end{equation}
For the source basis, the PBW order places \(Y\) after every factor of
\(X\) and before the inducing subalgebra. These complements are used
only as vector spaces; they need not be Lie subalgebras.

Write \(\Delta=\ad F\). The finite commutation formula
\eqref{eq:finite-commutation} gives
\[
\begin{aligned}
 h_iF^{-m}&=F^{-m}h_i+2mF^{-m-1}f_{i+b},\\
 e_iF^{-m}&=F^{-m}e_i
       -mF^{-m-1}(h_{i+b}+i\delta_{i+b,0}K)\\
 &\hspace{2em}-m(m+1)F^{-m-2}f_{i+2b}.
\end{aligned}
\]
It follows that
\[
 h_iw_1=(\lambda_i+2\delta_{i,0})w_1\quad(i\ge0),
 \qquad Kw_1=cw_1:
\]
if \(i>0\), then \(f_{i+b}u=0\), whereas for \(i=0\) the correction
term is \(2F^{-1}u\). Also, \(f_jw_1=0\) for \(j>b\), and at the
boundary \(f_bw_1=u+M=0\). If \(i\ge a+2\), then
\(i+b\ge L+1>0\), whence
\[
 e_iu=h_{i+b}u=f_{i+2b}u=0,\qquad\delta_{i+b,0}=0.
\]
Thus \(e_iw_1=0\). This verifies every inducing relation for
\(S_{a+1,b-1}\).

Extend \(\chi^+\) to the algebra character
\(\chi_U^+:U(S_{a+1,b-1})\to\C\). By taking finite products and linear
combinations of the verified relations, we obtain
\(sw_1=\chi_U^+(s)w_1\) for every \(s\in U(S_{a+1,b-1})\). Hence the
linear map
\[
 U\otimes_\C\C_{\chi^+}\longrightarrow T_FM,\qquad
 x\otimes\zeta1_{\chi^+}\longmapsto\zeta xw_1
\]
annihilates every balancing relation, since
\[
 xs\otimes\zeta1_{\chi^+}-x\otimes s(\zeta1_{\chi^+})
 \longmapsto
 \zeta x\bigl(sw_1-\chi_U^+(s)w_1\bigr)=0.
\]
It therefore descends to a well-defined map \(\Phi:A\to T_FM\).
The map commutes with left multiplication by every element of \(U\),
so it is a module homomorphism. Uniqueness follows from \(A=Uv\).
No use of \(\lambda_L\ne0\) has been made at this stage.

\subsection{Explicit preimages of pure negative powers}

Since \(Yu=0\), \([Y,F]=h_L\), and \(f_{b+L}u=0\), with no central
term because \(L>0\), the preceding formulas yield
\[
 Yw_m=-m\lambda_Lw_{m+1},\qquad
 Y^pw_1=(-1)^pp!\lambda_L^pw_{p+1}.
\]
Therefore
\begin{equation}\label{eq:pure-preimage}
 q_m=\frac{(-1)^{m-1}}{(m-1)!\lambda_L^{m-1}}Y^{m-1}v
 \qquad(m\ge1)
\end{equation}
satisfies \(\Phi(q_m)=w_m\). In particular, every \(f_b^{-m}u+M\)
already has an explicit preimage. This calculation does not assume
surjectivity.

\subsection{The surjectivity induction and injectivity}

Assign weights to the generators of \(U\) by
\[
 \wt(e_i)=2,\qquad\wt(h_i)=1,\qquad\wt(f_i)=\wt(K)=0.
\]
Let \(U_{\le r}\) be the span of words of total weight at most \(r\),
with \(U_{\le r}=0\) for \(r<0\). Every nonzero bracket of generators
has weight at most the sum of their weights minus one, so PBW
reordering does not increase weight. On each generator, \(\Delta\)
lowers weight by at least one whenever its value is nonzero. The
Leibniz rule therefore gives
\[
 \Delta^j(X)\in U_{\le\rho(X)-j},\qquad
 \Delta^j(X)=0\quad(j>\rho(X)).
\]
Reorder \(\Delta^j(X)u\) in the PBW basis of \(M\). Inducing factors
act by scalars or zero, which again cannot increase weight. Thus
\[
 \Delta^j(X)u=\sum_{q,X'}c_{j,q,X'}F^qX'u,\qquad
 q\ge0,\quad
 c_{j,q,X'}\ne0\ \Longrightarrow\ \rho(X')\le\rho(X)-j,
\]
where the sum is finite. Substitute this into
\eqref{eq:finite-commutation}. The \(j=0\) term of \(Xw_m\) is
exactly \(B_{m,X}\). For \(j\ge1\), a term with \(q\ge m+j\) vanishes
in the quotient; otherwise it is a multiple of \(B_{m+j-q,X'}\)
with \(\rho(X')<\rho(X)\). Hence
\begin{equation}\label{eq:triangular}
 Xw_m=B_{m,X}
       +\text{a finite sum of terms of strictly smaller }\rho.
\end{equation}

We prove by induction on \(r\ge0\) that \(B_{m,X}\) lies in
\(\operatorname{im}\Phi\) for \emph{every} \(m\ge1\) and
\(\rho(X)\le r\).

\emph{Base case \(r=0\).}
Then \(X\) contains only \(f\)-factors and commutes with \(F\). Thus
\[
 B_{m,X}=Xw_m=\Phi(Xq_m).
\]
This includes every \(f\)-monomial, not only \(X=1\).

\emph{Induction step.}
Assume the assertion for \(r-1\). Fix any \(m\ge1\) and any \(X\)
with \(\rho(X)=r\). By \eqref{eq:triangular},
\[
 \Phi(Xq_m)=B_{m,X}+\sum_{\nu=1}^{s}c_\nu B_{m_\nu,X_\nu},
 \qquad m_\nu\ge1,\quad\rho(X_\nu)\le r-1.
\]
For each \(\nu\), the induction hypothesis supplies \(a_\nu\in A\)
with \(\Phi(a_\nu)=B_{m_\nu,X_\nu}\). Therefore
\[
 B_{m,X}=\Phi\left(Xq_m-\sum_{\nu=1}^{s}c_\nu a_\nu\right).
\]
Each step involves only finitely many terms, and the nonnegative
integer \(\rho\) strictly decreases in every correction term, so the
recursion terminates. Crucially, the induction covers all denominator
orders simultaneously; an increase in the denominator order causes
no difficulty. By the target basis in \eqref{eq:pbw-bases}, \(\Phi\)
is surjective.

For injectivity, \eqref{eq:triangular} implies
\[
 \Phi(XY^pv)=(-1)^pp!\lambda_L^pB_{p+1,X}
       +\text{terms of strictly smaller }\rho.
\]
In a nonzero finite linear combination of source basis vectors,
consider the largest occurring value of \(\rho\). Its terms map to
distinct target basis vectors, with coefficients multiplied by the
nonzero scalars \((-1)^pp!\lambda_L^p\). They cannot cancel. The image
is therefore nonzero, proving injectivity and completing the proof of
Theorem~\ref{thm:boundary}. Neither finite-dimensional filtration
pieces nor simplicity are needed.

\subsection{Special cases, iteration, and the fixed-module obstruction}

Taking \(a=-1\) and \(b=n\ge1\) gives the direct generalization
\begin{equation}\label{eq:fn-special}
 M_{0,n-1}(\chi^+)\cong T_{f_n}M_{-1,n}(\chi),
 \qquad\lambda_n\ne0,
\end{equation}
with
\[
 q_m=\frac{(-1)^{m-1}}{(m-1)!\lambda_n^{m-1}}e_0^{m-1}v.
\]
For successive boundary moves, put
\(N_0=M_{a,b}(\chi)\) and \(N_{r+1}=T_{f_{b-r}}N_r\). Then, for every
finite \(r\ge0\),
\[
 N_r\cong M_{a+r,b-r}(\chi^{(r)}),\qquad
 \chi^{(r)}(h_0)=\lambda_0+2r,
\]
with all other character values unchanged. Indeed, \(a+b\), \(L\),
and the nonzero value \(\lambda_L\) are preserved at each step, so
the theorem applies repeatedly. This is a statement about finitely
many successive localization quotients of changing modules, not about
infinite iteration or simultaneous localization.

The situation differs if the original module
\(M=M_{-1,1}(\chi)\) is kept fixed.

\emph{If \(n\ge2\),} then \(F=f_n\) annihilates \(u\). For any
\(Xu\), choose \(N\ge1\) with \((\ad F)^N(X)=0\). In the expansion
of \(F^NXu\), each term with \(j<N\) contains a positive power of
\(F\) acting on \(u\), while the term with \(j=N\) is zero.
Thus \(F\) acts locally nilpotently, and \(\D_FM=T_FM=0\).

\emph{If \(n\le0\),} then \(F=f_n\) belongs to the PBW complement.
Placing it first shows that it acts injectively and that the quotient
is nonzero. However,
\[
 h_1(F^{-1}u+M)
 =\lambda_1(F^{-1}u+M)+2(F^{-2}f_{n+1}u+M).
\]
Since \(n+1\le1\) and \(n+1\ne n\), the correction term is a nonzero
PBW basis vector, linearly independent of \(F^{-1}u+M\). Hence
\(F^{-1}u+M\) is not an \(h_1\)-eigenvector and cannot be the image
of the canonical inducing vector of a standard module of the above
form. This rules out only the specified cyclic map: it neither
classifies the quotient nor excludes an isomorphism using a different
cyclic vector.

Finally, if \(\lambda_L=0\), the cyclic map in
Theorem~\ref{thm:boundary} remains well-defined, but
\(\Phi(Yv)=0\) whereas \(Yv\ne0\). Thus that map is not an isomorphism.

\section{Adjoining the derivation: the correct tensor action}

For any \(\g\)-module \(M\), define
\[
 \I(M)=\widetilde U\otimes_U M.
\]
Ordering \(d\) before a basis of \(\g\) gives the right \(U\)-module
decomposition
\[
 \widetilde U=\bigoplus_{k\ge0}d^kU,\qquad
 \I(M)\cong\C[z]\otimes M,\qquad
 P(z)\otimes m\longleftrightarrow P(d)\otimes_U m.
\]
Here \(z\) is a formal polynomial variable, distinguished from the
derivation \(d\). Flipping the vector-space factors gives the proposed
space \(M\otimes\C[d]\); we keep the polynomial factor on the left.
The PBW decomposition also shows that \(\I\) is exact: kernels and
images after induction are computed coefficientwise.

\begin{lemma}\label{lem:shift-action}
Under this identification, the induced action is
\begin{equation}\label{eq:shift-action}
\begin{aligned}
 d(P(z)\otimes m)&=zP(z)\otimes m,\\
 x_r(P(z)\otimes m)&=P(z-r)\otimes x_rm,\\
 K(P(z)\otimes m)&=P(z)\otimes Km.
\end{aligned}
\end{equation}
These formulas define a \(\gt\)-module and realize \(\I(M)\).
\end{lemma}
\begin{proof}
The identity \(x_rd=(d-r)x_r\) gives
\(x_rP(d)=P(d-r)x_r\) by induction, so the formulas follow from the
induced action. They may also be checked directly. For homogeneous
modes,
\[
 [x_r,y_s](P(z)\otimes m)=P(z-r-s)\otimes[x_r,y_s]m.
\]
The noncentral bracket term has degree \(r+s\), and a central term
can occur only when \(r+s=0\), exactly as required. Moreover,
\[
\begin{aligned}
 [d,x_r](P(z)\otimes m)
 &=\bigl[zP(z-r)-(z-r)P(z-r)\bigr]\otimes x_rm\\
 &=rP(z-r)\otimes x_rm.
\end{aligned}
\]
Together with centrality of \(K\), this verifies all relations.
\end{proof}

The unshifted action \(x_r(P\otimes m)=P\otimes x_rm\) would give
\([d,x_r]=0\), contradicting the required relation whenever a
nonzero-degree mode acts nontrivially. Nor can one simply extend the
one-dimensional inducing character to include \(d\): any such
extension, denoted \(\widehat\chi\), would satisfy
\[
 0=[\widehat\chi(d),\widehat\chi(h_L)]
   =\widehat\chi([d,h_L])=L\lambda_L\ne0.
\]
Thus we retain \(S_{a,b}\), without \(d\), as the inducing subalgebra.
Transitivity of induction gives
\begin{equation}\label{eq:induction-transitivity}
 \widetilde M_{a,b}(\chi)
 :=\widetilde U\otimes_{U(S_{a,b})}\C_\chi
 \cong\I(M_{a,b}(\chi)).
\end{equation}

\section{Localization commutes with adjoining the derivation}

This section does not require \(M\) to be a particular induced module.
Fix \(n\in\Z\) and \(F=f_n\). For \(\Delta=\ad F\), the brackets give
\[
\begin{gathered}
 \Delta(f_r)=\Delta(K)=0,\qquad
 \Delta(h_r)=2f_{r+n},\qquad \Delta(d)=-nF,\\
 \Delta(e_r)=-(h_{r+n}+r\delta_{r+n,0}K),\qquad
 \Delta^2(e_r)=-2f_{r+2n},\qquad \Delta^3(e_r)=0.
\end{gathered}
\]
The Leibniz rule implies that \(\Delta\) is locally nilpotent on
both \(U\) and \(\widetilde U\). Their PBW associated graded algebras
are domains, so \(F\) is regular on both sides in each algebra.
If \(\Delta^{t+1}(s)=0\), then for \(N\ge t\),
\[
 F^Ns=\sum_{j=0}^t\binom Nj\Delta^j(s)F^{N-j},\qquad
 sF^N=\sum_{j=0}^t(-1)^j\binom NjF^{N-j}\Delta^j(s).
\]
Taking \(N\ge m+t\) gives the required right or left factor \(F^m\),
proving both Ore conditions. Module localization can be represented
by fractions, with
\[
 F^{-m}v=F^{-\ell}w
 \quad\Longleftrightarrow\quad
 F^{k-m}v=F^{k-\ell}w
 \quad\text{for some }k\ge\max\{m,\ell\}.
\]
Thus finite sums admit a common denominator, and the kernel of
\(M\to\D_FM\) consists exactly of vectors annihilated by some power
of \(F\).

\begin{theorem}\label{thm:commute}
Suppose \(F=f_n\) acts injectively on the \(\g\)-module \(M\).
There are natural \(\gt\)-module isomorphisms
\[
 \alpha_M:\D_F\I(M)\xrightarrow{\sim}\I(\D_FM),\qquad
 \overline\alpha_M:T_F\I(M)\xrightarrow{\sim}\I(T_FM).
\]
The first isomorphism and its inverse are given by
\begin{equation}\label{eq:alpha}
\begin{aligned}
 \alpha_M\bigl(F^{-m}(P(z)\otimes u)\bigr)
   &=P(z+mn)\otimes F^{-m}u,\\
 \alpha_M^{-1}\bigl(P(z)\otimes F^{-m}u\bigr)
   &=F^{-m}(P(z-mn)\otimes u),
\end{aligned}
\end{equation}
where \(m\ge0\), \(P\in\C[z]\), and \(u\in M\). The second
isomorphism is obtained by passing to quotients.
\end{theorem}

\begin{proof}
\emph{Injective action and construction.}
On \(\I(M)\), the action of \(F\) is
\[
 P(z)\otimes u\longmapsto P(z-n)\otimes Fu.
\]
Polynomial translation is invertible and \(F\) is injective on \(M\);
hence \(F\) is injective on \(\I(M)\). On \(\I(\D_FM)\), it is
invertible, with
\[
 F^{-1}(P(z)\otimes v)=P(z+n)\otimes F^{-1}v.
\]
Iteration gives the positive shift \(+mn\) in \eqref{eq:alpha}.
The basic fraction relation is respected because
\[
 P(z-n+(m+1)n)\otimes F^{-m-1}Fu
 =P(z+mn)\otimes F^{-m}u.
\]
Together with common denominators, this proves that \(\alpha_M\)
is well-defined.

\emph{Module compatibility.}
Local nilpotence gives the finite commutation identity
\begin{equation}\label{eq:finite-commutation}
 sF^{-m}
 =\sum_{j\ge0}\binom{m+j-1}{j}F^{-m-j}\Delta^j(s)
 \qquad(m\ge1).
\end{equation}
For \(m=1\), repeatedly substitute
\(sF^{-1}=F^{-1}s+F^{-1}\Delta(s)F^{-1}\).
The general case follows by induction using the binomial summation
identity obtained from Pascal's identity. Each nonzero
\(\Delta^j(x_r)\) is homogeneous of degree \(r+jn\); a possible
central term has degree zero. Therefore, for \(m\ge1\),
\begin{align*}
 &\alpha_M\bigl(x_rF^{-m}(P(z)\otimes u)\bigr)\\
 &\quad=\sum_{j\ge0}\binom{m+j-1}{j}
 P\bigl(z-(r+jn)+(m+j)n\bigr)
       \otimes F^{-m-j}\Delta^j(x_r)u\\
 &\quad=P(z+mn-r)\otimes x_rF^{-m}u\\
 &\quad=x_r\alpha_M\bigl(F^{-m}(P(z)\otimes u)\bigr).
\end{align*}
For \(m=0\), the assertion is immediate without commuting a
denominator. Compatibility with \(K\) follows directly. For the
derivation, \([d,F^{-m}]=-mnF^{-m}\), so
\begin{align*}
 \alpha_M\bigl(dF^{-m}(P(z)\otimes u)\bigr)
 &=\alpha_M\bigl(F^{-m}((z-mn)P(z)\otimes u)\bigr)\\
 &=zP(z+mn)\otimes F^{-m}u\\
 &=d\alpha_M\bigl(F^{-m}(P(z)\otimes u)\bigr).
\end{align*}

\emph{Bijectivity and passage to quotients.}
The second formula in \eqref{eq:alpha} supplies a preimage for every
elementary tensor, proving surjectivity. If
\(\alpha_M(F^{-m}w)=0\), apply \(F^m\) in the target: the image of
\(w\) in \(\I(\D_FM)\) is zero. Since \(M\hookrightarrow\D_FM\)
is injective and \(\I\) is exact,
\(\I(M)\hookrightarrow\I(\D_FM)\) is injective. Thus \(w=0\),
proving injectivity and the stated inverse formula.

On elements with denominator exponent zero, \(\alpha_M\) is the
induced inclusion. Applying \(\I\) to
\[
 0\longrightarrow M\longrightarrow\D_FM
 \longrightarrow T_FM\longrightarrow0
\]
therefore gives the quotient isomorphism. Finally, for a module map
\(g:M\to N\), the localized and induced maps send
\[
 F^{-m}u\longmapsto F^{-m}g(u),\qquad
 P\otimes u\longmapsto P\otimes g(u).
\]
Substitution in \eqref{eq:alpha} proves naturality.
\end{proof}

\begin{corollary}\label{cor:lift}
Under the hypotheses of Theorem~\ref{thm:commute}, a \(\g\)-module
isomorphism \(\Phi:A\xrightarrow{\sim}T_FM\) lifts to
\[
 \widetilde\Phi=\overline\alpha_M^{-1}\circ\I(\Phi):
 \I(A)\xrightarrow{\sim}T_F\I(M).
\]
If \(\Phi(a)=F^{-m}u+M\), then
\begin{equation}\label{eq:lift}
 \widetilde\Phi(P(z)\otimes a)
 =F^{-m}(P(z-mn)\otimes u)+\I(M).
\end{equation}
In particular, if \(\Phi(q_m)=F^{-m}u+M\), then
\begin{equation}\label{eq:poly-preimage}
 \boxed{\ 
 \widetilde\Phi^{-1}\bigl(F^{-m}(P(z)\otimes u)+\I(M)\bigr)
 =P(z+mn)\otimes q_m.\ }
\end{equation}
\end{corollary}

The induced map \(\I(\Phi)\) itself is indeed
\(\mathrm{id}_{\C[z]}\otimes\Phi\). However, identifying its target
with \(T_F\I(M)\) requires the translation in \eqref{eq:alpha}.
When \(n\ne0\), replacing \(P\otimes(F^{-m}u+M)\) by
\(F^{-m}(P\otimes u)+\I(M)\) without this translation is generally
incompatible with \(d\). When \(n=0\), the translations disappear.

\section{The combined isomorphism}

Apply Corollary~\ref{cor:lift} to Theorem~\ref{thm:boundary} and use
\eqref{eq:induction-transitivity}.

\begin{corollary}\label{cor:final}
If \(a+b\ge0\), \(L=a+b+1\), and \(\lambda_L\ne0\), then for every
central parameter \(c\),
\begin{equation}\label{eq:full-iso}
 \boxed{\ 
 \widetilde M_{a+1,b-1}(\chi^+)
 \xrightarrow[\ \widetilde\Phi\ ]{\ \sim\ }
 T_{f_b}\widetilde M_{a,b}(\chi).\ }
\end{equation}
The map sends \(\widetilde v=1\otimes v\) to
\(f_b^{-1}\widetilde u+\widetilde M_{a,b}(\chi)\), where
\(\widetilde u=1\otimes u\). For \(q_m\) from
\eqref{eq:pure-preimage},
\begin{equation}\label{eq:full-preimage}
\begin{aligned}
 &\widetilde\Phi^{-1}
 \bigl(f_b^{-m}(P(z)\otimes u)+\widetilde M_{a,b}(\chi)\bigr)\\
 &\hspace{5em}=P(z+mb)\otimes q_m.
\end{aligned}
\end{equation}
In particular, for \(a=-1\) and \(b=n\ge1\),
\[
 \widetilde M_{0,n-1}(\chi^+)
 \cong T_{f_n}\widetilde M_{-1,n}(\chi),
 \qquad\lambda_n\ne0.
\]
\end{corollary}

This completes the proposed procedure. Well-definedness, the
surjectivity induction, and injectivity first establish \(\Phi\)
without \(d\). Exact induction gives \(\I(\Phi)\), and
\(\overline\alpha_M^{-1}\) identifies its target with the localization
quotient after adjoining \(d\). No new surjectivity induction on the
polynomial degree is needed. Likewise, each finite sequence of
boundary moves lifts step by step.

\section{Smoothness and limitations}

If \(M\) is smooth, the finitely many coefficients \(m_i\) in
\(\sum_iP_i(z)\otimes m_i\) have a common bound above which all modes
annihilate them. By \eqref{eq:shift-action}, the same bound works
for this tensor. Thus \(\I(M)\) is smooth.

For an element \(F^{-m}u\) of the localization, with \(m\ge1\),
the finite commutation formula gives
\[
\begin{aligned}
 h_rF^{-m}u&=F^{-m}h_ru+2mF^{-m-1}f_{r+n}u,\\
 e_rF^{-m}u&=F^{-m}e_ru
       -mF^{-m-1}(h_{r+n}+r\delta_{r+n,0}K)u\\
 &\hspace{2em}-m(m+1)F^{-m-2}f_{r+2n}u,
\end{aligned}
\]
while \(f_r\) commutes with \(F\). Choose \(r\) sufficiently large
that \(r,r+n,r+2n\) exceed the annihilation bound for \(u\) and
\(r+n\ne0\). All terms then vanish. Common denominators and finite
sums show that \(\D_FM\), and hence \(T_FM\), are smooth. This
argument also covers negative integers \(n\).

The modules \(M_{a,b}(\chi)\) themselves are smooth: commuting a
sufficiently high mode through a fixed finite PBW word introduces
only finitely many index shifts, and the resulting high modes
annihilate the inducing vector.

An isomorphism statement does not by itself imply simplicity. In
general, \(\I\) does not preserve simple modules: the trivial
one-dimensional \(\g\)-module \(\C\) is simple, but
\(\I(\C)=\C[z]\) has the nonzero proper submodule \(z\C[z]\),
with all loop modes acting as zero and \(d\) acting by multiplication
by \(z\). This example only shows that induction does not preserve
simplicity in general; it does not satisfy the injectivity
hypothesis on \(F\).

For the noncritical induced modules considered here,
\cite[\S4.1]{FGXZ} uses an additional Casimir central reduction to
construct simple modules. The isomorphisms in this note alone do
not establish simplicity of the entire induced module with \(d\).

\appendix
\section{Scope and record of the Danus trial}

The trial used the public Codex release of Danus,
\texttt{v0.1.0-codex}, through a bounded native Windows adaptation.
It retained the upstream deterministic prechecks, verification entry
point, candidate-submission gate, and fact graph, while replacing
the Linux daemon launch mechanism with local read-only Codex calls.
No unsandboxed launch script was run, no additional paid strategic
consultation was enabled, and no global agent configuration was
changed.

Separate workers submitted candidates for the boundary isomorphism
and for compatibility with the derivation. Each review used an
independent cold-start verifier. The boundary candidate was accepted
on its first submission, with fact identifier
\texttt{d6bf6fef62dc8971}. The derivation candidate was initially
rejected because its theorem statement did not fully define the
ambient algebras and functor. After these definitions were supplied,
it was accepted with fact identifier \texttt{b388fabb701946d3}.
The second review confirmed consistency of the positive shift
\(+mn\), inverse shift \(-mn\), and derivation correction.

The two English proof candidates in the fact graph retain their
submitted statements, proofs, and verification results. The present
note is an edited English version of the integrated Chinese note:
it uses unified notation and combines the results functorially in
Corollary~\ref{cor:final}. Neither integrated version has undergone
a separate whole-paper verifier review. All 67 offline tests for
the adaptation and relevant upstream modules passed; dependency
deprecation warnings and nonblocking heuristic symbol-recognition
warnings were also recorded. These software tests are not
mathematical proofs.

The Danus verifier is an LLM, not Lean or Coq; acceptance does not
constitute a formal proof certificate \cite{DanusTrust}. The arguments
remain subject to author and advisor review. A complete classification
of nonboundary quotients of the fixed original module, and further
simplicity questions, lie outside the results established here.

\begin{thebibliography}{9}
\bibitem{FGXZ}
V.~Futorny, X.~Guo, Y.~Xue and K.~Zhao,
\emph{Smooth representations of affine Kac--Moody algebras},
arXiv:2404.03855v2 (2024), \S2.3 and \S4.
\url{https://arxiv.org/abs/2404.03855v2}.
\bibitem{DanusTrust}
FrenzyMath, \emph{Danus: Security and Trust Model},
\texttt{v0.1.0-codex}, \S1--4.
\url{https://github.com/frenzymath/Danus/blob/v0.1.0-codex/docs/security-and-trust.md}.
\end{thebibliography}
\end{document}
